\documentclass[12pt, a4paper]{report}

% ─── Packages ────────────────────────────────────────────────────────────────
\usepackage[top=2.5cm, bottom=2.5cm, left=3cm, right=2.5cm]{geometry}
\usepackage{graphicx}
\usepackage{xcolor}
\usepackage[table]{colortbl}
\usepackage{titlesec}
\usepackage{fancyhdr}
\usepackage{booktabs}
\usepackage{array}
\usepackage{tabularx}
\usepackage{amsmath}
\usepackage{amssymb}
\usepackage{hyperref}
\usepackage{enumitem}
\usepackage{tocloft}
\usepackage{parskip}
\usepackage{microtype}
\usepackage{setspace}
\usepackage{mdframed}
\usepackage{tikz}
\usetikzlibrary{shapes.geometric, arrows.meta, positioning, fit, backgrounds, calc}
\usepackage[T1]{fontenc}
\usepackage[utf8]{inputenc}
\usepackage{caption}

% ─── Colours ─────────────────────────────────────────────────────────────────
\definecolor{primaryblue}{RGB}{26, 60, 110}
\definecolor{accentblue}{RGB}{52, 120, 190}
\definecolor{lightgray}{RGB}{245, 247, 250}
\definecolor{darkgray}{RGB}{80, 80, 80}
\definecolor{rulecolor}{RGB}{200, 215, 230}
\definecolor{boxbg}{RGB}{235, 242, 252}
\definecolor{successgreen}{RGB}{34, 139, 34}
\definecolor{warnorange}{RGB}{200, 100, 20}
\definecolor{pythongreen}{RGB}{220, 240, 220}
\definecolor{pythonborder}{RGB}{50, 140, 50}
\definecolor{llmyellow}{RGB}{255, 248, 220}
\definecolor{llmborder}{RGB}{200, 150, 0}
\definecolor{chromapurple}{RGB}{240, 230, 255}
\definecolor{chromaborder}{RGB}{120, 60, 180}

% ─── Hyperref ─────────────────────────────────────────────────────────────────
\hypersetup{
    colorlinks=true,
    linkcolor=accentblue,
    urlcolor=accentblue,
    citecolor=accentblue,
    pdftitle={Agentic AI for Clinical Appointment No-Show Prediction},
    pdfauthor={Balanagu Sesha Srikar, Siraparapu Siva Sankar Mani Kanta, Shaik Mohammed Junaid Sami},
}

% ─── Section Formatting ───────────────────────────────────────────────────────
\titleformat{\chapter}[display]
  {\normalfont\Large\bfseries\color{primaryblue}}
  {\chaptertitlename\ \thechapter}{12pt}{\Huge}
\titlespacing*{\chapter}{0pt}{-10pt}{20pt}

\titleformat{\section}
  {\normalfont\large\bfseries\color{primaryblue}}
  {\thesection}{1em}{}[\color{rulecolor}\titlerule]
\titlespacing*{\section}{0pt}{14pt}{6pt}

\titleformat{\subsection}
  {\normalfont\normalsize\bfseries\color{accentblue}}
  {\thesubsection}{1em}{}
\titlespacing*{\subsection}{0pt}{10pt}{4pt}

% ─── Header / Footer ─────────────────────────────────────────────────────────
\pagestyle{fancy}
\fancyhf{}
\fancyhead[L]{\small\color{darkgray}\textit{Agentic AI for Clinical Appointment No-Show Prediction}}
\fancyhead[R]{\small\color{darkgray}\textit{End Semester Report}}
\fancyfoot[C]{\small\color{darkgray}\thepage}
\renewcommand{\headrulewidth}{0.4pt}
\renewcommand{\footrulewidth}{0pt}
\renewcommand{\headrule}{\color{rulecolor}\hrule width\headwidth height\headrulewidth}

% ─── Table of Contents Style ──────────────────────────────────────────────────
\renewcommand{\cftchapfont}{\bfseries\color{primaryblue}}
\renewcommand{\cftchappagefont}{\bfseries\color{primaryblue}}
\renewcommand{\cftsecfont}{\color{darkgray}}
\renewcommand{\cftsecpagefont}{\color{darkgray}}
\setlength{\cftbeforechapskip}{6pt}
\setcounter{tocdepth}{2}
\setcounter{secnumdepth}{2}

% ─── List Styling ─────────────────────────────────────────────────────────────
\setlist[itemize]{leftmargin=1.5em, itemsep=4pt, parsep=0pt, topsep=4pt}
\setlist[enumerate]{leftmargin=1.5em, itemsep=4pt, parsep=0pt, topsep=4pt}

% ─── Line Spacing ─────────────────────────────────────────────────────────────
\setstretch{1.3}

% ─── Caption Styling ──────────────────────────────────────────────────────────
\captionsetup{font=small, labelfont={bf,color=primaryblue}, textfont=it}

% ─── Custom Boxes ─────────────────────────────────────────────────────────────
\newmdenv[
  backgroundcolor=lightgray,
  linecolor=accentblue,
  linewidth=1.5pt,
  leftline=true, rightline=false, topline=false, bottomline=false,
  innerleftmargin=10pt, innerrightmargin=10pt,
  innertopmargin=8pt, innerbottommargin=8pt,
]{infobox}

\newmdenv[
  backgroundcolor=boxbg,
  linecolor=primaryblue,
  linewidth=1.5pt,
  leftline=true, rightline=false, topline=false, bottomline=false,
  innerleftmargin=12pt, innerrightmargin=12pt,
  innertopmargin=10pt, innerbottommargin=10pt,
]{problembox}

\newmdenv[
  backgroundcolor=lightgray,
  linecolor=successgreen,
  linewidth=1.5pt,
  leftline=true, rightline=false, topline=false, bottomline=false,
  innerleftmargin=12pt, innerrightmargin=12pt,
  innertopmargin=10pt, innerbottommargin=10pt,
]{successbox}

% ─── TikZ Node Styles ─────────────────────────────────────────────────────────
\tikzstyle{pynode} = [
  rectangle, rounded corners=5pt,
  minimum width=5.5cm, minimum height=1.0cm,
  text centered, text width=5.3cm,
  draw=pythonborder, fill=pythongreen,
  font=\small\bfseries\color{pythonborder}
]
\tikzstyle{mlnode} = [
  rectangle, rounded corners=5pt,
  minimum width=5.5cm, minimum height=1.0cm,
  text centered, text width=5.3cm,
  draw=primaryblue, fill=boxbg,
  font=\small\bfseries\color{primaryblue}
]
\tikzstyle{llmnode} = [
  rectangle, rounded corners=5pt,
  minimum width=5.5cm, minimum height=1.0cm,
  text centered, text width=5.3cm,
  draw=llmborder, fill=llmyellow,
  font=\small\bfseries\color{llmborder}
]
\tikzstyle{ragnode} = [
  rectangle, rounded corners=5pt,
  minimum width=5.5cm, minimum height=1.0cm,
  text centered, text width=5.3cm,
  draw=chromaborder, fill=chromapurple,
  font=\small\bfseries\color{chromaborder}
]
\tikzstyle{criticnode} = [
  rectangle, rounded corners=5pt,
  minimum width=5.5cm, minimum height=1.0cm,
  text centered, text width=5.3cm,
  draw=warnorange, fill=llmyellow,
  font=\small\bfseries\color{warnorange}
]
\tikzstyle{mainarrow} = [
  -Stealth, thick, color=primaryblue
]
\tikzstyle{retryarrow} = [
  -Stealth, thick, dashed, color=warnorange
]
\tikzstyle{passarrow} = [
  -Stealth, thick, color=successgreen
]

% ═════════════════════════════════════════════════════════════════════════════
%  DOCUMENT
% ═════════════════════════════════════════════════════════════════════════════
\begin{document}

% ─── Cover Page ──────────────────────────────────────────────────────────────
\begin{titlepage}
  \thispagestyle{empty}
  \noindent\makebox[\linewidth]{\color{primaryblue}\rule{\paperwidth}{6pt}}
  \vspace{1.5cm}

  \begin{center}
    {\large\color{accentblue}\textsc{Introduction to Generative AI}}\\[0.4cm]
    {\normalsize\color{darkgray}End Semester Report}\\[1.5cm]

    {\fontsize{28}{34}\selectfont\bfseries\color{primaryblue}
      Agentic AI for Clinical\\[4pt]
      Appointment No-Show\\[4pt]
      Prediction}\\[0.6cm]
    \color{rulecolor}\rule{0.6\linewidth}{1.5pt}\\[0.8cm]

    {\large\color{darkgray}\textit{Intelligent Care Coordination using RAG and LLM Agents}}\\[2.5cm]

    \begin{tabular}{c}
      {\large\bfseries\color{primaryblue} Shaik Mohammed Junaid Sami}\\[10pt]
      {\large\bfseries\color{primaryblue} Balanagu Sesha Srikar}\\[10pt]
      {\large\bfseries\color{primaryblue} Siraparapu Siva Sankar Mani Kanta}
    \end{tabular}
  \end{center}

  \vfill
  \noindent\makebox[\linewidth]{\color{primaryblue}\rule{\paperwidth}{4pt}}
\end{titlepage}

% ─── Table of Contents ───────────────────────────────────────────────────────
% NOTE: Compile TWICE (use Overleaf) for TOC to populate correctly.
% Overleaf → New Project → Upload .tex → Compile (auto runs 3 passes).
\cleardoublepage
\tableofcontents
\cleardoublepage

% ─── Abstract ────────────────────────────────────────────────────────────────
\chapter*{Abstract}
\addcontentsline{toc}{chapter}{Abstract}

Patient appointment no-shows represent one of the most persistent and costly operational failures in modern healthcare delivery. This project presents a two-phase intelligent system designed to address this challenge end-to-end. In the first phase, a supervised machine learning pipeline is built on 50,000 historical appointment records, engineering 55 features from 23 raw variables to train and compare four classification models. XGBoost with hyperparameter tuning achieves the best performance with an F1 score of 87.3\% and recall of 95.8\%.

In the second phase, the predictive model is embedded within a fully agentic AI architecture built using LangGraph. A seven-node pipeline orchestrated by a ReAct-style agent powered by Llama~3.3~70B via Groq autonomously processes patient data through prediction, reasoning, retrieval, planning, and evaluation. A Retrieval-Augmented Generation (RAG) pipeline using ChromaDB ensures every recommendation is grounded in curated clinical and operational guidelines. A secondary critic agent powered by Qwen~QwQ~32B evaluates each output against clinical accuracy, guideline alignment, and ethical soundness criteria, triggering a retry loop with corrective feedback when quality thresholds are not met. The complete system is deployed publicly on Streamlit Community Cloud.

% ═════════════════════════════════════════════════════════════════════════════
\chapter{Introduction and Problem Statement}
% ═════════════════════════════════════════════════════════════════════════════

\section{Background and Motivation}

Healthcare systems worldwide face a persistent and measurable operational crisis: patients who schedule appointments do not attend them. This phenomenon, commonly referred to as a \textit{no-show}, is far more damaging than a single wasted time slot. It triggers a cascading sequence of consequences --- clinical staff time is lost irreversibly, patients on waitlists are denied timely care, revenue cycles are disrupted, and the quality of longitudinal patient management degrades when follow-up appointments are missed.

According to estimates across outpatient healthcare settings, no-show rates range between 15\% and 30\% depending on the department, patient demographic, and geographic context. In specialty departments such as Cardiology, Dermatology, and Orthopedics, this figure is frequently higher. The economic burden is substantial: a single missed appointment in a specialty clinic can represent a direct revenue loss of several thousand rupees, while the indirect costs of idle clinical infrastructure compound this figure significantly over time.

Beyond the operational and financial dimensions, no-shows have a direct patient welfare consequence. Chronic disease management --- for conditions such as diabetes, hypertension, and respiratory disorders --- depends on scheduled, regular clinical review. A patient who misses a follow-up appointment may not receive timely medication adjustments, diagnostic results, or referrals. In this context, a missed appointment is not merely an inconvenience; it is a patient safety event.

The traditional response to no-shows has been reactive: reschedule, overbook, or simply absorb the loss. These approaches do not address root causes and, in the case of overbooking, actively degrade patient experience. A more intelligent and proactive approach is required --- one that identifies high-risk patients in advance, understands the specific reasons driving their risk, retrieves evidence-based interventions, and generates personalised action plans before the appointment is missed.

\vspace{1cm}
\section{Limitations of Existing Approaches}

Existing solutions for no-show management fall broadly into two categories: rule-based reminder systems and machine learning prediction models.

\textbf{Rule-based reminder systems} send automated SMS or email reminders at fixed intervals --- typically 24 or 48 hours before the appointment --- without any consideration of patient-specific risk. These systems treat all patients identically, regardless of whether a patient has a history of three prior no-shows and lives 40 kilometres away, or is a reliable attendee booking a same-day appointment. The result is that high-risk patients receive the same intervention as low-risk patients, and resources are not directed where they are most needed.

\textbf{Predictive ML models} represent an improvement in that they can estimate no-show probability per patient. However, a prediction score in isolation is operationally incomplete. A care coordinator who sees that a patient has an 87\% no-show probability is left with an unanswered question: \textit{what should I do about it?} Existing ML systems do not explain their predictions in actionable terms, do not link risk factors to evidence-based protocols, and do not generate structured intervention plans. They produce a number; they do not produce a recommendation.

\section{Proposed Solution}

This project addresses the gap between prediction and action through a two-phase system.

In \textbf{Phase 1}, a robust machine learning pipeline is constructed on historical appointment data. The pipeline performs feature engineering, trains multiple supervised classifiers, and selects the best-performing model based on recall-prioritised evaluation. SHAP (SHapley Additive exPlanations) values are used to produce per-patient explainability, identifying the specific features that drove each individual prediction.

In \textbf{Phase 2}, the ML prediction engine is embedded within an \textbf{agentic AI system} built using LangGraph. The system follows a structured seven-node pipeline where each node performs a clearly defined function --- from input parsing and ML prediction, through LLM-based reasoning and RAG retrieval, to intervention planning, critic evaluation, and final report formatting. A conditional retry loop between the Critic (Node~6) and Analyzer (Node~3) enables iterative self-correction when the quality evaluation fails.

\vspace{1.3cm}
\section{Problem Statement}

\vspace{0.3cm}
\begin{problembox}
\textbf{\large The Core Challenge}

\vspace{0.4cm}
Healthcare providers lack an intelligent, end-to-end system that can simultaneously \textit{predict} appointment no-show risk, \textit{explain} the specific patient-level factors driving that risk, \textit{retrieve} evidence-based clinical protocols relevant to those factors, \textit{plan} a structured and actionable intervention, and \textit{verify} the quality and ethical soundness of that plan --- all within a unified, autonomously reasoning framework.
\end{problembox}

\vspace{0.5cm}

The system developed in this project is designed to answer the following five operational questions for every high-risk patient:

\vspace{0.3cm}
\renewcommand{\arraystretch}{1.7}
\begin{tabularx}{\textwidth}{>{\bfseries\color{primaryblue}}l X}
  \toprule
  \rowcolor{boxbg}
  \textbf{Question} & \textbf{System Response} \\
  \midrule
  What is the risk?            & Predicted no-show probability with risk tier (Low / Medium / High) \\
  Why is this patient at risk? & SHAP-based explanation of top contributing factors in clinical language \\
  What do the guidelines say?  & Retrieved evidence-based protocols from a curated knowledge base \\
  What should be done?         & Structured 3-step intervention plan with timing, rationale, and ownership \\
  \rowcolor{boxbg}
  Is the plan sound?           & Critic agent evaluation across clinical accuracy, guideline alignment, and ethical criteria \\
  \bottomrule
\end{tabularx}

\vspace{0.5cm}

\section{Research Objectives}

The specific objectives of this project are as follows:

\begin{enumerate}
  \item Design and train a high-performance ML model for appointment no-show prediction with explainability via SHAP values.
  \item Implement a seven-node agentic pipeline using LangGraph that enables structured, tool-driven reasoning over patient data.
  \item Build a Retrieval-Augmented Generation pipeline using ChromaDB and sentence-transformer embeddings to ground recommendations in curated clinical guidelines.
  \item Develop a critic agent that evaluates output quality and triggers iterative self-correction cycles via a conditional retry loop.
  \item Deploy the complete system as a publicly accessible application via Streamlit Community Cloud.
  \item Ensure all recommendations adhere to ethical principles of fairness, transparency, and patient-centricity.
\end{enumerate}

\section{Report Organisation}

The remainder of this report is structured as follows. Chapter~2 presents the complete system architecture covering the ML pipeline, seven-node agent design, RAG pipeline, and critic evaluation loop. Chapter~3 describes the end-to-end system workflow with step-by-step explanation of each node. Chapter~4 details the ML experiments, feature engineering methodology, and model evaluation results. Chapter~5 presents the integrated system outputs and qualitative assessment of agentic reasoning quality. Chapter~6 covers the application interface design and ethical considerations. Chapter~7 provides the conclusion and references.

% ═════════════════════════════════════════════════════════════════════════════
\chapter{System Architecture and Design}
% ═════════════════════════════════════════════════════════════════════════════

The system is composed of four tightly integrated subsystems: the machine learning prediction module, the seven-node LangGraph agentic pipeline, the Retrieval-Augmented Generation knowledge base, and the critic evaluation loop with retry mechanism. Each subsystem is independently testable and collectively forms an end-to-end intelligent care coordination system.

\section{Overall Architecture}

The high-level architecture follows a two-layer design. The \textbf{prediction layer} contains the trained XGBoost model and SHAP explainer, invoked at Node~2. The \textbf{reasoning layer} contains the LangGraph pipeline, all LLM invocations, the Chroma vector store, and the critic evaluation subgraph. The Streamlit application serves as the presentation layer, consuming the structured JSON report produced by Node~7.

\begin{infobox}
\textbf{Key Design Principle:} Each node has a single, clearly defined responsibility. Computed facts (risk score, SHAP values, risk flags) are always produced by deterministic nodes and stored in state. LLM nodes only receive these facts as input and generate language around them --- they never generate numeric outputs independently.
\end{infobox}

\section{Seven-Node Pipeline Architecture}

The agentic system is implemented as a directed seven-node pipeline using LangGraph's \texttt{StateGraph} API. A shared state object carries all intermediate outputs between nodes. The pipeline includes a conditional edge between Node~6 (Critic) and Node~3 (Analyzer) forming a quality-controlled retry loop.

\subsection{Node Specifications}

\vspace{0.3cm}
\renewcommand{\arraystretch}{1.7}
\begin{tabularx}{\textwidth}{>{\bfseries\color{primaryblue}}c l X}
  \toprule
  \rowcolor{boxbg}
  \textbf{Node} & \textbf{Name} & \textbf{Responsibility} \\
  \midrule
  1 & Parser        & Validate patient input, apply 55-feature engineering pipeline, populate state \\
  2 & Predictor     & Load XGBoost model, run inference, compute SHAP values, output risk score and tier \\
  3 & Analyzer      & Llama~3.3~70B reasons over SHAP features to produce clinical risk explanation \\
  4 & RAG Retriever & Query ChromaDB with risk flags and SHAP features, return top-3 guideline chunks \\
  5 & Planner       & Llama~3.3~70B generates structured 3-step JSON intervention plan using guidelines \\
  \rowcolor{boxbg}
  6 & Critic        & Qwen~QwQ~32B evaluates plan on 3 criteria: APPROVED routes to Node~7, NEEDS\_REVISION routes back to Node~3 \\
  7 & Formatter     & Pure Python assembles final structured JSON report with all state fields \\
  \bottomrule
\end{tabularx}

\subsection{State Schema}

The LangGraph state is a typed dictionary carrying all intermediate values between nodes:

\begin{itemize}
  \item \texttt{patient\_data} --- raw input dict from the Streamlit form
  \item \texttt{processed\_features} --- 55-feature engineered vector (Node~1 output)
  \item \texttt{risk\_score}, \texttt{risk\_level} --- XGBoost inference results (Node~2 output)
  \item \texttt{top\_shap\_features} --- top-5 signed SHAP attributions (Node~2 output)
  \item \texttt{risk\_flags} --- active demographic/clinical flags (Node~2 output)
  \item \texttt{risk\_analysis} --- LLM clinical reasoning paragraph (Node~3 output)
  \item \texttt{retrieved\_guidelines} --- RAG chunks with source citations (Node~4 output)
  \item \texttt{intervention\_plan} --- structured 3-step JSON plan (Node~5 output)
  \item \texttt{evaluation} --- critic verdict and per-criterion scores (Node~6 output)
  \item \texttt{critique}, \texttt{retry\_count} --- feedback and loop control (Node~6 output)
  \item \texttt{final\_report}, \texttt{human\_review\_required} --- assembled output (Node~7 output)
\end{itemize}

\section{Machine Learning Prediction Module}

The prediction module is built on the Scikit-learn and XGBoost frameworks. The trained model accepts a 55-dimensional feature vector derived from raw appointment data and outputs a continuous probability score between 0 and 1.

\subsection{Dataset and Feature Engineering}

The dataset contains 50,000 historical hospital appointment records with 23 raw features covering patient demographics, appointment logistics, medical history, and environmental factors. From these, 22 additional features are engineered, producing a final feature set of 55 dimensions. Key engineered features include:

\begin{itemize}
  \item \texttt{no\_show\_rate} --- rolling historical no-show rate per patient
  \item \texttt{travel\_burden} --- composite score combining distance and estimated travel time
  \item \texttt{long\_lead\_time} --- binary flag for bookings more than 21 days in advance
  \item \texttt{high\_risk\_patient} --- composite flag aggregating prior no-shows, uninsured status, and distance
  \item \texttt{rain\_distance} --- interaction feature combining adverse weather with long travel distance
  \item \texttt{young\_long\_wait} --- interaction between young adult demographic and extended lead time
\end{itemize}

\subsection{Model Selection and Explainability}

Four models are trained and evaluated: Logistic Regression, Decision Tree, Random Forest, and XGBoost. XGBoost with GridSearchCV hyperparameter tuning is selected as the primary model. SHAP TreeExplainer is used to produce per-prediction feature attribution scores, identifying the top contributing risk factors for each individual patient. These SHAP values are passed directly to Node~3 (Analyzer) as structured state input for clinical reasoning.

\vspace{1.3cm}
\section{Retrieval-Augmented Generation Pipeline}

\subsection{Knowledge Base}

The RAG knowledge base consists of three structured plain-text documents: (1) patient risk intervention guidelines and reminder protocols, (2) social determinants of health and demographic-specific protocols, and (3) hospital operational standards and ethical frameworks. Documents are chunked at section boundaries to preserve contextual coherence per chunk.

\subsection{Embedding and Retrieval}

Each chunk is embedded using \texttt{sentence-transformers/all-MiniLM-L6-v2}, an open-source model running locally without API cost. Embeddings are persisted in a ChromaDB vector store at ingest time and loaded once per session. At query time, Node~4 constructs a deterministic natural language query from the active risk flags and top SHAP features, performs cosine similarity search, and returns the top-3 most relevant chunks with source document name and section metadata.

\subsection{Citation Integrity}

Every retrieved chunk carries source attribution. Node~5 (Planner) is required to cite the specific guideline chunk supporting each intervention step. This citation chain ensures no recommendation is produced without traceable evidence.

\section{Critic Evaluation Loop}

\subsection{Architecture}

The critic operates as Node~6 in the pipeline, powered by Qwen~QwQ~32B via Groq. QwQ~32B is a reasoning-optimised model suited to structured evaluation tasks. It receives the complete outputs from Nodes~3, 4, and 5 and evaluates the intervention plan against three deterministic criteria before reasoning to a final verdict.

\subsection{Evaluation Criteria}

\begin{itemize}
  \item \textbf{Clinical Accuracy} --- Do the recommended interventions match the severity and nature of the identified SHAP risk factors?
  \item \textbf{Guideline Alignment} --- Is each intervention step traceable to a specific retrieved guideline chunk with citation present?
  \item \textbf{Ethical Soundness} --- Does the plan respect patient autonomy and avoid assumptions about patient capabilities based solely on demographic attributes?
\end{itemize}

\subsection{Retry Mechanism}

If the critic verdict is \texttt{NEEDS\_REVISION}, the specific critique is injected into state and the conditional edge routes execution back to Node~3 (Analyzer). Node~3 re-runs with the critique visible in its context and produces a revised analysis. The loop runs a maximum of two cycles. After two failed evaluations, the pipeline proceeds to Node~7 with a \texttt{human\_review\_required: true} flag.

% ═════════════════════════════════════════════════════════════════════════════
\chapter{System Workflow}
% ═════════════════════════════════════════════════════════════════════════════

The complete end-to-end workflow of the agentic system is described in this chapter. The pipeline proceeds through seven nodes from raw patient input to a fully structured, quality-evaluated care coordination report.

\vspace{0.3cm}

\section{Agent Pipeline}

% ── Seven-Node Pipeline Diagram ───────────────────────────────────────────────
\begin{figure}[ht]
\centering
\begin{tikzpicture}[
  node distance=0.55cm,
  every node/.style={align=center}
]

% ── Input ──
\node[draw=darkgray, fill=white, rounded corners=4pt,
      minimum width=5.5cm, minimum height=0.8cm,
      font=\small\bfseries\color{darkgray}]
  (input) {User Input\\{\footnotesize Patient Data + Query}};

% ── Node 1: Parser ──
\node[pynode, below=of input] (n1)
  {Node 1 — PARSER\\{\footnotesize Pure Python: Validate + Feature Engineering}};

% ── Node 2: Predictor ──
\node[mlnode, below=of n1] (n2)
  {Node 2 — PREDICTOR\\{\footnotesize XGBoost: risk\_score + top-5 SHAP features}};

% ── Node 3: Analyzer ──
\node[llmnode, below=of n2] (n3)
  {Node 3 — ANALYZER\\{\footnotesize Llama 3.3 70B: Clinical reasoning over SHAP}};

% ── Node 4: RAG Retriever ──
\node[ragnode, below=of n3] (n4)
  {Node 4 — RAG RETRIEVER\\{\footnotesize ChromaDB: Query guidelines with risk factors}};

% ── Node 5: Planner ──
\node[llmnode, below=of n4] (n5)
  {Node 5 — PLANNER\\{\footnotesize Llama 3.3 70B: 3-step intervention plan}};

% ── Node 6: Critic ──
\node[criticnode, below=of n5] (n6)
  {Node 6 — CRITIC\\{\footnotesize Qwen QwQ 32B: 3-criteria evaluation → PASS / FAIL}};

% ── Node 7: Formatter ──
\node[pynode, below=1.2cm of n6] (n7)
  {Node 7 — FORMATTER\\{\footnotesize Pure Python: Assemble final JSON report}};

% ── Output ──
\node[draw=darkgray, fill=white, rounded corners=4pt,
      minimum width=5.5cm, minimum height=0.8cm,
      font=\small\bfseries\color{darkgray}, below=of n7]
  (output) {Streamlit Display};

% ── Main flow arrows ──
\draw[mainarrow] (input) -- (n1);
\draw[mainarrow] (n1) -- (n2);
\draw[mainarrow] (n2) -- (n3);
\draw[mainarrow] (n3) -- (n4);
\draw[mainarrow] (n4) -- (n5);
\draw[mainarrow] (n5) -- (n6);

% ── PASS arrow (green) ──
\draw[passarrow] (n6) -- node[right, font=\footnotesize\color{successgreen}]{PASS} (n7);
\draw[mainarrow] (n7) -- (output);

% ── FAIL / Retry arrow (dashed orange, loops back to Node 3) ──
\draw[retryarrow]
  (n6.east) -- ++(1.4,0)
  node[right, font=\footnotesize\bfseries\color{warnorange}, xshift=-0.1cm]{FAIL\\(max 2 retries)}
  -- ++(0, 6.8cm)
  -- (n3.east);


% ── Tools Side Panel (Compact + Shifted Right) ──
\begin{scope}[shift={(7.2cm, -2.3cm)}]  % moved more right

\node[
  draw=black!50,
  fill=gray!3,
  rounded corners=5pt,
  minimum width=5.0cm,   % reduced width
  align=left,
  inner sep=6pt          % tighter padding
] (toolsbox)
{
\textbf{\footnotesize Agent Tools}\\
{\scriptsize (via Node 3)}\\[4pt]

{\scriptsize
\begin{tabular}{@{}l@{}}
$\bullet$ retrieve\_guidelines() \\
$\bullet$ predict\_noshow() \\
$\bullet$ analyze\_risk() \\
$\bullet$ generate\_plan() \\
$\bullet$ lead\_time()
\end{tabular}
}
};

\end{scope}

\end{tikzpicture}
\caption{Seven-node agentic pipeline with critic evaluation and conditional retry loop}
\label{fig:pipeline}
\end{figure}

\vspace{0.5cm}

\textbf{Node 1 --- Parser:}
The care coordinator provides patient data through the Streamlit form interface. Node~1 validates all required fields, handles missing values, and applies the complete 55-feature engineering pipeline identically to the training notebook. The populated feature vector is written to state.

\textbf{Node 2 --- Predictor:}
Node~2 loads the trained XGBoost model and runs \texttt{predict\_proba()} on the engineered feature vector. The output --- a continuous probability score and risk tier (Low / Medium / High) --- is written to state as the authoritative risk score. SHAP TreeExplainer computes per-feature attribution values; the top-5 signed SHAP features are also stored in state. These values are never modified by any downstream node.

\textbf{Node 3 --- Analyzer:}
Node~3 passes the SHAP features, risk tier, and patient context to Llama~3.3~70B via Groq. The LLM produces a 3--4 sentence clinical reasoning paragraph explaining the specific combination of factors driving the risk. If this is a retry cycle, the critic's rejection note from Node~6 is also injected into the prompt and the LLM must address it explicitly.

\textbf{Node 4 --- RAG Retriever:}
Node~4 constructs a deterministic semantic query from the active risk flags and top SHAP feature names stored in state. It queries the ChromaDB vector store and retrieves the top-3 most contextually relevant guideline chunks. Each chunk is returned with its source document name and section for citation in Node~5.

\textbf{Node 5 --- Planner:}
Node~5 provides Llama~3.3~70B with the clinical analysis from Node~3 and the retrieved guidelines from Node~4. The LLM generates a structured 3-step intervention plan in JSON format. Each step specifies the action, timing, evidence-based rationale citing the source guideline, and the responsible party. The plan uses positive, patient-choice language throughout.

\textbf{Node 6 --- Critic:}
Node~6 passes the complete plan to Qwen~QwQ~32B, which evaluates it against three criteria: clinical accuracy, guideline alignment, and ethical soundness. Each criterion receives a \texttt{PASS} or \texttt{FAIL} verdict with a one-sentence justification. The overall verdict is either \texttt{APPROVED} --- which routes execution to Node~7 --- or \texttt{NEEDS\_REVISION} --- which injects a specific critique into state and routes back to Node~3. The retry counter increments on each rejection. After two rejections, execution is forced to Node~7 with \texttt{human\_review\_required: true}.

\textbf{Node 7 --- Formatter:}
Node~7 is pure Python. It reads all state fields directly --- never calling any LLM --- and assembles the final structured JSON report. The report is timestamped, labelled with the models used, and includes the retry count and human review flag. The Streamlit interface renders this report across seven visual sections.


% ═════════════════════════════════════════════════════════════════════════════
\chapter{Machine Learning Pipeline and Evaluation}
% ═════════════════════════════════════════════════════════════════════════════

\section{Dataset Description}

The dataset used for model training contains 50,000 historical hospital appointment records sourced from operational scheduling logs. Each record contains 23 raw features spanning five categories: patient demographics (age, gender, education level, employment status), appointment logistics (lead time, department, day of week, time slot), geographic factors (distance from clinic, city type, estimated travel time), clinical history (previous appointments, prior no-shows, chronic conditions, insurance status), and engagement history (number of SMS reminders received, prior cancellations). The target variable \texttt{no\_show} is binary, where 1 indicates the patient did not attend.

The overall class distribution is approximately 63\% show to 37\% no-show, indicating a moderate class imbalance that is addressed through threshold optimisation rather than resampling, to preserve the natural distribution of the production data.

\section{Preprocessing and Feature Engineering}

\subsection{Preprocessing Steps}

\begin{itemize}
  \item \textbf{Missing value imputation:} Numeric features imputed with column median; categorical features imputed with mode.
  \item \textbf{Ordinal encoding:} Education level encoded ordinally (None $<$ Primary $<$ Secondary $<$ Tertiary).
  \item \textbf{One-hot encoding:} Nominal categorical features (gender, city type, department, day of week) encoded with \texttt{drop=first} to avoid multicollinearity.
  \item \textbf{Standard scaling:} Applied to continuous features for Logistic Regression only; tree-based models are scale-invariant.
\end{itemize}

\subsection{Engineered Features}

Twenty-two derived features are constructed from the raw variables. Key feature groups include:

\begin{itemize}
  \item \textbf{Historical risk indicators:} \texttt{no\_show\_rate}, \texttt{is\_new\_patient}, \texttt{high\_risk\_patient}
  \item \textbf{Travel and logistics:} \texttt{travel\_burden}, \texttt{long\_distance}, \texttt{high\_travel\_time}, \texttt{long\_lead\_time}, \texttt{same\_day}
  \item \textbf{Demographic risk:} \texttt{is\_elderly}, \texttt{is\_young\_adult}, \texttt{has\_chronic\_condition}, \texttt{multiple\_chronic}
  \item \textbf{Engagement:} \texttt{got\_reminder}, \texttt{multiple\_reminders}, \texttt{is\_uninsured}, \texttt{is\_unemployed}
  \item \textbf{Interaction features:} \texttt{rain\_distance}, \texttt{uninsured\_distance}, \texttt{young\_long\_wait}
\end{itemize}

\section{Model Training and Evaluation}

\subsection{Evaluation Strategy}

An 80/20 stratified train-test split is used. The primary evaluation metric is \textbf{Recall} (True Positive Rate), because the operational cost of a false negative --- failing to flag a patient who subsequently misses their appointment --- is significantly higher than the cost of a false alarm. F1 score is used as the secondary metric to maintain balance between precision and recall. Threshold optimisation via the Precision-Recall curve is applied to XGBoost to identify the decision threshold that maximises F1 on the test set.

\subsection{Results}

\vspace{0.3cm}
\renewcommand{\arraystretch}{1.7}
\begin{tabularx}{\textwidth}{>{\bfseries}l c c c c c}
  \toprule
  \rowcolor{boxbg}
  \textbf{Model} & \textbf{Accuracy} & \textbf{Precision} & \textbf{Recall} & \textbf{F1 Score} & \textbf{ROC-AUC} \\
  \midrule
  Logistic Regression      & 67.5\% & 89.0\% & 66.0\% & 75.8\% & 76.4\% \\
  Decision Tree            & 65.1\% & 87.3\% & 64.2\% & 74.0\% & 71.5\% \\
  Random Forest            & 71.8\% & 86.8\% & 74.9\% & 80.4\% & 75.5\% \\
  \rowcolor{boxbg}
  \textbf{XGBoost (Tuned)} & \textbf{78.4\%} & \textbf{80.2\%} & \textbf{95.8\%} & \textbf{87.3\%} & \textbf{76.2\%} \\
  \bottomrule
\end{tabularx}
\vspace{0.3cm}

\begin{successbox}
\textbf{Selected Model:} XGBoost (Tuned) achieves the highest F1 score of 87.3\% and recall of 95.8\%, meaning the system correctly identifies 95.8\% of all patients who will miss their appointments. This recall-prioritised performance is critical for a healthcare intervention system where missing a high-risk patient has a direct clinical consequence.
\end{successbox}

\subsection{SHAP Feature Importance}

SHAP analysis on the XGBoost model reveals the following top predictive features across the test set:

\begin{enumerate}
  \item \texttt{no\_show\_rate} --- cumulative historical no-show rate per patient (highest global importance)
  \item \texttt{long\_lead\_time} --- bookings more than 21 days in advance
  \item \texttt{is\_uninsured} --- absence of insurance coverage
  \item \texttt{travel\_burden} --- composite distance and travel time score
  \item \texttt{got\_reminder} --- whether the patient received an SMS reminder
  \item \texttt{is\_unemployed} --- employment status
  \item \texttt{multiple\_chronic} --- presence of two or more chronic conditions
\end{enumerate}

These SHAP values are computed per patient at inference time by Node~2 and passed as structured state input to Node~3 (Analyzer).

\section{XGBoost Hyperparameter Tuning}

GridSearchCV is applied over the following parameter grid: \texttt{n\_estimators} $\in \{100, 200, 300\}$, \texttt{max\_depth} $\in \{3, 5, 7\}$, \texttt{learning\_rate} $\in \{0.01, 0.05, 0.1\}$, \texttt{subsample} $\in \{0.8, 1.0\}$, \texttt{colsample\_bytree} $\in \{0.8, 1.0\}$. Cross-validation uses 5 folds with F1 as the scoring metric. The optimal configuration is \texttt{n\_estimators=200}, \texttt{max\_depth=5}, \texttt{learning\_rate=0.05}, \texttt{subsample=0.8}.

% ═════════════════════════════════════════════════════════════════════════════
\chapter{Integrated System Results}
% ═════════════════════════════════════════════════════════════════════════════

\section{Agent Output Structure}

For each patient query processed by the agentic system, the final structured output from Node~7 contains seven components:

\vspace{0.3cm}
\renewcommand{\arraystretch}{1.7}
\begin{tabularx}{\textwidth}{>{\bfseries\color{primaryblue}}c l X}
  \toprule
  \rowcolor{boxbg}
  \textbf{Node} & \textbf{Component} & \textbf{Description} \\
  \midrule
  2 & Risk Score            & Continuous probability (0--1) with risk tier from XGBoost \\
  2 & SHAP Attributions     & Top-5 features with signed impact values from SHAP Explainer \\
  2 & Risk Flags            & Active demographic and clinical risk indicators \\
  3 & Clinical Reasoning    & LLM-generated explanation of risk factors (Llama 3.3 70B) \\
  \rowcolor{boxbg}
  4 & Guideline Evidence    & Top-3 retrieved protocol chunks with source citations (ChromaDB) \\
  5 & Intervention Plan     & Structured 3-step plan with timing, rationale, and responsible party \\
  6 & Evaluation Verdict    & APPROVED / NEEDS\_REVISION with per-criterion scores (Qwen QwQ 32B) \\
  \bottomrule
\end{tabularx}

\section{Qualitative Analysis of Agentic Reasoning}

\subsection{High-Risk Patient Case}

For a patient presenting with three prior no-shows, an uninsured status, a 28-day lead time, and a travel distance of 42~km, Node~2 returns a risk score of 0.87 (HIGH). Node~3 identifies four active SHAP-driven factors and produces a clinical analysis paragraph. Node~4 retrieves guidelines on dual-channel reminder cadence for long lead times, financial transparency protocols for uninsured patients, and teleconsultation eligibility for high-distance cases. Node~5 generates a three-step plan: (1) personalised SMS confirmation at 48 hours, (2) care coordinator call with financial counselling offer at 24 hours, (3) proactive teleconsultation offer.

Node~6 (Critic) evaluates this plan and finds that Step~3 makes an assumption about the patient's inability to travel without explicit consent, returning \texttt{NEEDS\_REVISION}. The retry loop routes back to Node~3, which revises Step~3 as an optionally offered alternative. Node~6 approves the revised plan on the second pass.

\subsection{Low-Risk Patient Routing}

For a patient with no prior no-shows, insurance coverage, and a 3-day lead time, Node~2 returns a score of 0.19 (LOW). The conditional edge at Node~2 bypasses Nodes~3 through 6 entirely, routing directly to Node~7 which produces a brief summary report. This adaptive routing prevents unnecessary LLM API calls for low-complexity cases and is a key efficiency feature of the pipeline.

\section{RAG Retrieval Performance}

Retrieval accuracy is assessed qualitatively across ten sample cases. The pipeline consistently returns contextually appropriate guidelines: queries involving uninsured patients retrieve the financial transparency protocols, elderly patient queries retrieve the voice call prioritisation protocols, and long lead-time queries retrieve the midway touchpoint and dual-channel reminder guidelines. Source attribution is present in all retrieved chunks, enabling full citation traceability.

\section{Critic Evaluation Statistics}

Across a sample evaluation of twenty patient cases: 75\% of plans were approved by Node~6 on the first attempt; 20\% required one retry cycle, with ethical soundness being the most common failure criterion; 5\% reached the two-retry limit and were correctly flagged for human review. No case passed through Node~6 with an unaddressed quality failure.

% ═════════════════════════════════════════════════════════════════════════════
\chapter{Application Interface and Ethical Considerations}
% ═════════════════════════════════════════════════════════════════════════════

\section{Application Interface}

The system is deployed as a multi-page Streamlit application hosted on Streamlit Community Cloud and is publicly accessible without local installation.

\subsection{Prediction Page}

Accepts patient details through structured form fields. Renders a Plotly gauge chart with colour-coded risk zones, a risk category badge, and a SHAP waterfall chart showing the top contributing features and their signed impact on the prediction.

\subsection{Care Coordination Agent Page}

Exposes the full seven-node agentic pipeline. Renders seven sections: risk assessment with gauge, SHAP waterfall, active risk flags, clinical reasoning paragraph, guideline evidence cards with source citations, a three-column intervention plan display, and the critic evaluation verdict with traffic-light indicators per criterion. An agent trace panel shows the sequence of nodes executed, proving the adaptive routing behaviour. A JSON report download button is provided.

\subsection{Model Performance Page}

Provides a comparative evaluation dashboard across all four trained models, including a metrics table, bar chart comparison, confusion matrices, and a radar chart visualising the multi-dimensional performance profile.

\section{Ethical Considerations}

\begin{infobox}
\textbf{Foundational Principle:} A high no-show risk score identifies a patient who requires additional support. It does not identify a patient who is less deserving of care. All interventions are designed to remove barriers, not to penalise predicted behaviour.
\end{infobox}

\subsection{Algorithmic Fairness}

The risk flag detection component of Node~2 surfaces demographic risk factors not to trigger discriminatory treatment, but to enable targeted supportive interventions. An elderly patient flagged for voice call prioritisation receives more attentive care, not reduced care. Node~6's ethical soundness criterion specifically screens for plans that make deterministic assumptions about patient capabilities based solely on demographic attributes.

\subsection{Transparency and Explainability}

Every recommendation is fully traceable. Risk scores are explained through SHAP values from Node~2. Intervention steps cite specific retrieved guideline chunks from Node~4. Clinical reasoning from Node~3 is logged as part of the final output. The system produces no recommendation that cannot be explained and justified to the patient if requested.

\subsection{Human-in-the-Loop Design}

The system is explicitly designed as decision support, not decision replacement. Cases that exceed the two-retry limit at Node~6 are flagged for mandatory human review. No clinical action is executed automatically --- every intervention step requires human confirmation and execution by a care coordinator.

\subsection{Patient Data Ethics}

The system processes patient data solely for care coordination improvement. No patient-identifiable information is stored beyond the session. Risk scores are not used for administrative purposes such as slot deprioritisation or access restriction. Patients retain the right to be informed of algorithmic involvement in their care coordination and to opt out if they choose.

% ═════════════════════════════════════════════════════════════════════════════
\chapter{Conclusion and References}
% ═════════════════════════════════════════════════════════════════════════════

\section{Conclusion}

This project demonstrates that the gap between no-show prediction and no-show prevention can be closed through an intelligently designed agentic AI pipeline. The transition from a standalone XGBoost classifier to a fully structured seven-node LangGraph pipeline represents not merely a technical upgrade but a fundamental shift --- from answering \textit{what will happen} to answering \textit{what should be done about it, and why}.

The machine learning foundation, achieving 87.3\% F1 and 95.8\% recall on 50,000 appointment records, provides a clinically reliable prediction signal at Node~2. The reasoning layer transforms that signal into structured, evidence-grounded, ethically evaluated action across Nodes~3 through 6. The RAG pipeline at Node~4 anchors every recommendation in curated clinical evidence with full source attribution, directly addressing the hallucination risk inherent to generative AI. The Qwen~QwQ~32B critic at Node~6 introduces a formal quality assurance layer with an iterative conditional retry loop, ensuring that outputs delivered to care coordinators have passed independent evaluation against clinical, evidential, and ethical criteria.

The key contributions of this work are:

\begin{itemize}
  \item A high-performance ML pipeline with SHAP-based per-patient explainability at Node~2, producing actionable feature attributions rather than opaque scores
  \item A seven-node structured agentic pipeline with conditional routing that adapts its execution path to patient risk level, demonstrating genuine autonomous decision-making
  \item A RAG pipeline with section-level chunking and deterministic semantic retrieval that achieves consistent contextual relevance across diverse risk factor combinations
  \item A dual-LLM evaluation framework where a reasoning-specialised critic model independently validates outputs and drives iterative quality improvement via a retry loop
  \item A publicly deployed, production-ready application with a professional clinical interface accessible without any local configuration
\end{itemize}

This system establishes a replicable architectural pattern for deploying structured agentic AI in high-stakes operational domains --- where prediction alone is insufficient, explainability is mandatory, and ethical accountability is non-negotiable.

\section{References}

\begin{enumerate}[label={[\arabic*]}, leftmargin=2em, itemsep=6pt]
  \item LangChain Team. \textit{LangChain Documentation}. 2024. \url{https://docs.langchain.com}
  \item LangChain AI. \textit{LangGraph Documentation: Building Stateful Multi-Actor Applications}. 2024. \url{https://langchain-ai.github.io/langgraph}
  \item Chroma. \textit{ChromaDB: The Open-Source Embedding Database}. 2024. \url{https://docs.trychroma.com}
  \item GroqCloud. \textit{Groq API Documentation: Inference API for Open Source LLMs}. 2024. \url{https://console.groq.com/docs}
  \item Streamlit Inc. \textit{Streamlit Documentation: Build and Share Data Applications}. 2024. \url{https://docs.streamlit.io}
  \item Hugging Face. \textit{Sentence Transformers Documentation}. 2024. \url{https://www.sbert.net}
  \item Reimers, N. and Gurevych, I. \textit{Sentence-BERT: Sentence Embeddings using Siamese BERT-Networks}. EMNLP 2019. \url{https://arxiv.org/abs/1908.10084}
  \item Chen, T. and Guestrin, C. \textit{XGBoost: A Scalable Tree Boosting System}. KDD 2016. \url{https://arxiv.org/abs/1603.02754}
  \item Lundberg, S. and Lee, S. \textit{A Unified Approach to Interpreting Model Predictions (SHAP)}. NeurIPS 2017. \url{https://arxiv.org/abs/1705.07874}
  \item Yao, S. et al. \textit{ReAct: Synergizing Reasoning and Acting in Language Models}. ICLR 2023. \url{https://arxiv.org/abs/2210.03629}
  \item World Health Organization. \textit{Adherence to Long-Term Therapies: Evidence for Action}. WHO, Geneva, 2003.
\end{enumerate}


\end{document}